% ============================================================
%  WNAR-SW: Weighted Negative Association Rule Mining with
%  Regularity in Medical Data Streams
%  Research Paper — LaTeX Source
% ============================================================
\documentclass[10pt,twocolumn]{article}

% ── Page geometry ────────────────────────────────────────────
\usepackage[
  top=1in, bottom=1in,
  left=0.75in, right=0.75in,
  columnsep=0.25in
]{geometry}

% ── Typography & encoding ─────────────────────────────────────
\usepackage[T1]{fontenc}
\usepackage[utf8]{inputenc}
\usepackage[protrusion=false,expansion=false]{microtype}

% ── Math ─────────────────────────────────────────────────────
\usepackage{amsmath,amssymb,amsthm,mathtools}

% ── Graphics & color ─────────────────────────────────────────
\usepackage{xcolor}
\usepackage{tikz}
\usepackage{pgfplots}
\pgfplotsset{compat=1.18}
\usetikzlibrary{arrows.meta, positioning, shapes.geometric, fit, backgrounds, calc}

% ── Tables ───────────────────────────────────────────────────
\usepackage{booktabs}
\usepackage{tabularx}
\usepackage{multirow}
\usepackage{array}

% ── Algorithms ───────────────────────────────────────────────
\usepackage[ruled,vlined,linesnumbered]{algorithm2e}
\SetKwComment{Comment}{$\triangleright$\ }{}

% ── Floats & captions ────────────────────────────────────────
\usepackage{float}
\usepackage{caption}
\captionsetup{font=small, labelfont=bf, skip=4pt}

% ── Lists ────────────────────────────────────────────────────
\usepackage{enumitem}
\setlist{noitemsep, topsep=3pt}

% ── Callout boxes ────────────────────────────────────────────
\usepackage{tcolorbox}
\tcbuselibrary{skins, breakable}

% ── Cross-references & hyperlinks ────────────────────────────
\usepackage{hyperref}
\usepackage{cleveref}
\hypersetup{
  colorlinks=true,
  linkcolor=darkblue,
  citecolor=darkblue,
  urlcolor=darkblue
}

% ── Section styling ──────────────────────────────────────────
\usepackage{titlesec}
\usepackage{abstract}

% ── Footnotes / misc ─────────────────────────────────────────
\usepackage{footmisc}
\usepackage{setspace}

% ── Colors ───────────────────────────────────────────────────
\definecolor{darkblue}{RGB}{23,55,94}
\definecolor{medblue}{RGB}{46,117,182}
\definecolor{lightblue}{RGB}{222,234,246}
\definecolor{lightgray}{RGB}{242,242,242}
\definecolor{accentgray}{RGB}{180,180,180}

% ── Section title formatting ──────────────────────────────────
\titleformat{\section}
  {\color{darkblue}\large\bfseries\scshape}
  {\thesection.}{0.5em}{}[{\color{medblue}\titlerule[0.5pt]}]

\titleformat{\subsection}
  {\color{medblue}\normalsize\bfseries}
  {\thesubsection.}{0.5em}{}

\titleformat{\subsubsection}
  {\color{darkblue}\normalsize\itshape\bfseries}
  {\thesubsubsection.}{0.5em}{}

\titlespacing*{\section}{0pt}{10pt}{4pt}
\titlespacing*{\subsection}{0pt}{7pt}{2pt}
\titlespacing*{\subsubsection}{0pt}{5pt}{1pt}

% ── Theorem environments ──────────────────────────────────────
\theoremstyle{plain}
\newtheorem{theorem}{Theorem}[section]
\newtheorem{lemma}[theorem]{Lemma}
\newtheorem{proposition}[theorem]{Proposition}
\newtheorem{corollary}[theorem]{Corollary}

\theoremstyle{definition}
\newtheorem{definition}{Definition}[section]
\newtheorem{example}{Example}[section]

\theoremstyle{remark}
\newtheorem{remark}{Remark}[section]
\newtheorem{property}{Property}[section]

% ── Callout tcolorbox styles ──────────────────────────────────
\tcbset{
  defbox/.style={
    enhanced, breakable,
    colback=lightblue!40, colframe=medblue,
    fonttitle=\bfseries\small, fontupper=\small,
    boxrule=0.5pt, arc=2pt,
    left=4pt, right=4pt, top=3pt, bottom=3pt,
    title style={colback=medblue, coltext=white}
  },
  notebox/.style={
    enhanced, breakable,
    colback=lightgray, colframe=accentgray,
    fontupper=\small\itshape,
    boxrule=0.4pt, arc=2pt,
    left=4pt, right=4pt, top=3pt, bottom=3pt,
    borderline west={2pt}{0pt}{medblue}
  }
}

% ── Custom commands ───────────────────────────────────────────
\newcommand{\ie}{\textit{i.e.},\xspace}
\newcommand{\eg}{\textit{e.g.},\xspace}
\newcommand{\etal}{\textit{et al.}\xspace}
\newcommand{\wsup}{\mathit{wsup}}
\newcommand{\wconf}{\mathit{wconf}}
\newcommand{\wNegSup}{\mathit{wNegSup}}
\newcommand{\wNegConf}{\mathit{wNegConf}}
\newcommand{\maxRegW}{\mathit{maxRegW}}
\newcommand{\minSup}{\mathit{minSup}}
\newcommand{\minConf}{\mathit{minConf}}
\newcommand{\SW}{\mathit{SW}}

% ── Abstract styling ──────────────────────────────────────────
\renewcommand{\abstractnamefont}{\bfseries\scshape\color{darkblue}}
\renewcommand{\abstracttextfont}{\small}
\setlength{\absleftindent}{0pt}
\setlength{\absrightindent}{0pt}

% ── Title block ───────────────────────────────────────────────
\title{%
  {\color{darkblue}\LARGE\bfseries
  WNAR-SW: Weighted Negative Association\\
  Rule Mining with Regularity in\\
  Medical Data Streams}\\[6pt]
  {\color{medblue}\normalsize\itshape
  Bridging the Gap Between Frequency-Regularity Mining\\
  and Weighted Association Rule Discovery}
}

\author{%
  \textbf{[Author Names]}\\
  \small Department of Computer Science\\
  \small [University Name]\\
  \small \texttt{[email@university.edu]}
}

\date{%
  \small\itshape Submitted to [Journal/Conference Name]\\
  \today
}

% ════════════════════════════════════════════════════════════
\begin{document}
% ════════════════════════════════════════════════════════════

\twocolumn[{%
  \maketitle
  \thispagestyle{empty}
  \begin{onecolabstract}
    \noindent
    Association rule mining is a cornerstone of data mining, yet the classical Apriori
    algorithm suffers from four compounding limitations when applied to real-time medical
    environments: it requires multiple database scans incompatible with streaming data;
    it discovers only positive associations while dangerous \emph{negative} drug interactions
    remain invisible; it treats all items as equally significant regardless of clinical
    severity; and it conflates sporadic coincidences with sustained, clinically meaningful
    patterns by ignoring temporal regularity.
    Two recent works address subsets of these limitations in isolation.
    Jammalamadaka and Budaraju~\cite{sastry2025} introduce a single-pass streaming
    algorithm that mines frequent, regular, and negative associations from medical data
    streams but applies no weighting to chemical compounds.
    Ouyang and Huang~\cite{ouyang2015} extend Apriori with weighted support over sliding
    windows to reflect item significance, yet mine only positive associations and apply
    no regularity constraint.
    This paper proposes \textbf{WNAR-SW} (Weighted Negative Association Rule Mining
    with Regularity in a Sliding Window), a unified algorithm that simultaneously
    addresses all four limitations.
    WNAR-SW introduces \emph{weighted regularity}, a severity-scaled temporal consistency
    constraint; \emph{weighted negative support}, which amplifies the signal of high-risk
    incompatible chemical pairs even when their raw frequency is low; and a
    \emph{severity-ranked alert} mechanism that triage physician notifications by
    clinical danger level.
    The algorithm degrades gracefully: setting all weights to unity reproduces the
    behaviour of~\cite{sastry2025}; disabling regularity yields a streaming
    weighted negative association miner.
    Theoretical analysis establishes the correctness of the weighted anti-monotonicity
    property under the negative mining extension.
    Experimental projections on a synthetic medical dataset of 100{,}000 records
    demonstrate that WNAR-SW eliminates the class of dangerous interactions that both
    prior algorithms miss while producing a ranked, actionable alert stream.
    \par\medskip
    \noindent
    \textbf{Keywords:} association rule mining, negative associations, weighted support,
    data streams, sliding window, regularity, medical data, drug interactions,
    adverse effects.
    \par\medskip
    \noindent
    \textbf{MSC (2020):} 68T05, 62P10, 68W27.
  \end{onecolabstract}
  \vspace{10pt}
  \hrule
  \vspace{8pt}
}]

% ─────────────────────────────────────────────────────────────
\section{Introduction}
\label{sec:intro}
% ─────────────────────────────────────────────────────────────

The problem of discovering hidden structure in large collections of
transactional records — known as \emph{association rule mining} — has
been studied intensively since the seminal work of Agrawal
\etal~\cite{agrawal1993,agrawal1994}.
The prototypical application, identifying which products are frequently
purchased together in a retail store, has since expanded to span
fraud detection, genomics, network security, and — most critically for
this work — clinical pharmacology.

The COVID-19 pandemic provided a stark reminder of what happens when
real-time drug interaction monitoring is absent.
Combinations such as beta-blockers with Remdesivir, or chemotherapy
agents with cardiac glycosides, were administered to patients without
immediate automated oversight.
Negative associations — the signal that two drugs or their constituent
chemicals systematically \emph{fail} to co-occur safely — were only
identified retrospectively, often after patient harm.
The challenge is not merely computational efficiency; it is the need
for a system that watches continuously flowing prescription data,
knows which chemicals are more dangerous than others, distinguishes
chronic patterns from coincidental spikes, and prioritises warnings
appropriately.

\subsection{Limitations of the Classical Approach}
\label{subsec:limitations}

The Apriori algorithm~\cite{agrawal1994}, the most widely taught and
deployed foundation for association rule mining, has four structural
limitations that collectively make it unsuitable for real-time medical
decision support:

\begin{enumerate}[label=\textbf{L\arabic*.}, leftmargin=*, itemsep=2pt]
  \item \textbf{Iterative scanning.} Apriori rescans the entire database
        once per itemset size level $k$. Streaming medical data — arriving
        continuously from wearables, electronic health records, and pharmacy
        dispensing systems — cannot be rescanned.

  \item \textbf{Positive-only rules.} Apriori finds rules of the form
        $X \Rightarrow Y$. Rules of the form $X \Rightarrow \lnot Y$ —
        signalling that $X$ and $Y$ cannot safely coexist — are
        structurally invisible to it.

  \item \textbf{Uniform item significance.} Every item, whether a
        harmless vitamin supplement or a potent chemotherapy agent,
        contributes identically to frequency counts.
        A rare but lethal incompatibility may be pruned before detection.

  \item \textbf{No temporal regularity.} Apriori counts occurrences
        but cannot distinguish a pattern that appears consistently every
        two days from one that appeared 100 times during a single outbreak
        and never again.
\end{enumerate}

\subsection{Contribution Summary}
\label{subsec:contributions}

This paper makes the following contributions:

\begin{itemize}[leftmargin=*]
  \item We formally define \emph{weighted regularity}
        $(\maxRegW)$, a severity-scaled temporal consistency threshold
        that tightens the regularity requirement for high-weight
        (high-danger) items (\cref{def:weighted_regularity}).

  \item We introduce \emph{weighted negative support}
        $(\wNegSup)$ and \emph{weighted negative confidence}
        $(\wNegConf)$ that jointly capture both clinical severity and
        the strength of mutual exclusion between chemical itemsets
        (\cref{def:wnegsup,def:wnegconf}).

  \item We propose the \textbf{WNAR-SW} algorithm, a three-process
        single-pass pipeline that mines negative associations from
        sliding-window medical data streams with weighted regularity
        filtering (\cref{sec:algorithm}).

  \item We prove that the weighted downward-closure property extends
        to the negative mining setting under WNAR-SW's weight
        formulation (\cref{thm:closure}).

  \item We define a \emph{severity score} that produces a ranked
        alert stream suitable for real-time clinical decision support
        (\cref{def:severity}).
\end{itemize}

\subsection{Paper Organisation}
\label{subsec:organisation}

\Cref{sec:related} surveys related work.
\Cref{sec:background} establishes formal preliminaries.
\Cref{sec:method} develops the WNAR-SW framework.
\Cref{sec:algorithm} presents the algorithm with pseudocode.
\Cref{sec:analysis} provides theoretical analysis.
\Cref{sec:discussion} discusses experimental projections and
practical considerations.
\Cref{sec:conclusion} concludes with directions for future work.

% ─────────────────────────────────────────────────────────────
\section{Related Work}
\label{sec:related}
% ─────────────────────────────────────────────────────────────

\subsection{Foundations of Association Rule Mining}

Association rule mining was introduced by Agrawal
\etal~\cite{agrawal1993}, with the Apriori algorithm formalised
in~\cite{agrawal1994}.
The core \emph{anti-monotonicity property} — that no superset of an
infrequent itemset can be frequent — enables exponential pruning of
the candidate space.
Zaki~\cite{zaki2003} demonstrated that vertical database
representations (Di-sets) accelerate itemset support counting through
bitwise intersection, a technique directly inherited by the papers
this work extends.

\subsection{Negative Association Rule Mining}

Savasere \etal~\cite{savasere1998} first studied strong negative
associations, requiring that item sets be individually frequent yet
mutually infrequent.
Antonie and Zaiane~\cite{antonie2004} proposed correlation-based
negative rule mining, generating negative rules only when two itemsets
are statistically negatively correlated.
Cornelis \etal~\cite{cornelis2006} systematically catalogued
algorithms for both positive and negative rules and identified the
critical gap that most methods fail to detect negative associations
among infrequent itemsets that carry high clinical weight.
Dong \etal~\cite{dong2006} introduced multi-confidence thresholds for
the four rule types ($A\!\Rightarrow\!B$, $A\!\Rightarrow\!\lnot B$,
$\lnot A\!\Rightarrow\!B$, $\lnot A\!\Rightarrow\!\lnot B$).

\subsection{Mining Association Rules from Data Streams}

Lin \etal~\cite{lin2006} developed a time-sensitive sliding window
approach for frequent itemset mining from streams.
Tanbeer \etal~\cite{tanbeer2010stream} proposed the Regular Pattern
Stream tree (RPS-tree) for detecting temporally regular patterns in
streaming data, introducing the regularity concept this work extends.
Tanbeer \etal~\cite{tanbeer2009sliding} further formalised
sliding-window frequent pattern mining.
Ouyang~\cite{ouyang2013} extended sliding-window methods to
simultaneously mine positive and negative association rules
(MPNAR-SW).

\subsection{Weighted Association Rule Mining}

Cai \etal~\cite{cai1998} introduced weighted association rules,
assigning significance scores to items to reflect real-world
importance.
Ouyang and Huang~\cite{ouyang2015} — the second foundational work
of this paper — extended weighted mining to data streams via
bit-sequence representations, proving that the downward-closure
property holds when weighting uses the maximum-weight item in each
itemset.
Bagui and Dhar~\cite{bagui2019} implemented positive and negative
association rule mining in a MapReduce framework, though without
weighting or regularity constraints.

\subsection{Medical Applications}

Jammalamadaka and Budaraju~\cite{sastry2025} — the first foundational
work of this paper — proposed a three-process streaming algorithm that
mines frequent, regular, negative associations specifically from
medical prescription data, demonstrating reduction of noise
associations from 0.73 to 0.43 per 1{,}000 itemsets through
combined frequency-regularity thresholds.
Wei \etal~\cite{wei2020} applied machine learning to predict adverse
drug response risk but did not consider real-time stream processing.
Zhang \etal~\cite{zhang2020} extracted drug-drug interactions from
biomedical literature but modelled only positive interactions.

\subsection{Gap Analysis}

\Cref{tab:comparison} summarises the feature coverage of prior work
against the requirements established in \cref{subsec:limitations}.
No prior algorithm satisfies all four requirements simultaneously.
WNAR-SW is designed to close this gap.

\begin{table}[H]
  \centering
  \caption{Feature coverage of related approaches.
  \checkmark$=$ supported, $\times=$ not supported.}
  \label{tab:comparison}
  \renewcommand{\arraystretch}{1.25}
  \small
  \begin{tabularx}{\columnwidth}{@{}lcccc@{}}
    \toprule
    \textbf{Method} & \textbf{L1} & \textbf{L2} & \textbf{L3} & \textbf{L4} \\
                    & Stream & Neg. & Weight & Reg. \\
    \midrule
    Apriori~\cite{agrawal1994}
      & $\times$ & $\times$ & $\times$ & $\times$ \\
    MPNAR-SW~\cite{ouyang2013}
      & \checkmark & \checkmark & $\times$ & $\times$ \\
    MWAR-SW~\cite{ouyang2015}
      & \checkmark & $\times$ & \checkmark & $\times$ \\
    Jammalamadaka~\cite{sastry2025}
      & \checkmark & \checkmark & $\times$ & \checkmark \\
    \textbf{WNAR-SW (ours)}
      & \checkmark & \checkmark & \checkmark & \checkmark \\
    \bottomrule
  \end{tabularx}
\end{table}

% ─────────────────────────────────────────────────────────────
\section{Background and Preliminaries}
\label{sec:background}
% ─────────────────────────────────────────────────────────────

\subsection{Standard Association Rules}

\begin{definition}[Transaction Database]
\label{def:database}
Let $\mathcal{I} = \{i_1, i_2, \ldots, i_m\}$ be a set of \emph{items}.
A \emph{transaction} is a pair $T = (\mathit{tid}, X)$ where
$\mathit{tid}$ is a unique identifier and $X \subseteq \mathcal{I}$.
A \emph{database} $\mathcal{D}$ is a finite multiset of transactions.
\end{definition}

\begin{definition}[Support and Confidence]
\label{def:support}
The \emph{support} of an itemset $X \subseteq \mathcal{I}$ in
$\mathcal{D}$ is
\[
  \mathit{sup}(X) = \frac{|\{T \in \mathcal{D} : X \subseteq T\}|}{|\mathcal{D}|}.
\]
An itemset $X$ is \emph{frequent} if $\mathit{sup}(X) \geq \minSup$
for a user-defined threshold $\minSup \in (0,1]$.
The \emph{confidence} of a rule $X \Rightarrow Y$ (where
$X \cap Y = \emptyset$) is
\[
  \mathit{conf}(X \Rightarrow Y)
    = \frac{\mathit{sup}(X \cup Y)}{\mathit{sup}(X)}.
\]
\end{definition}

\subsection{Negative Association Rules}

\begin{definition}[Negative Rule]
\label{def:negative_rule}
A \emph{negative association rule} $A \Rightarrow \lnot B$ holds
when itemsets $A$ and $B$ are each individually frequent but
jointly infrequent:
\[
  \mathit{sup}(A) \geq \minSup, \quad
  \mathit{sup}(B) \geq \minSup, \quad
  \mathit{sup}(A \cup B) < \minSup.
\]
Its support and confidence are:
\begin{align}
  \mathit{sup}(A \Rightarrow \lnot B)
    &= \mathit{sup}(A) - \mathit{sup}(A \cup B), \label{eq:neg_sup}\\
  \mathit{conf}(A \Rightarrow \lnot B)
    &= \frac{\mathit{sup}(A \Rightarrow \lnot B)}{\mathit{sup}(A)}. \label{eq:neg_conf}
\end{align}
\end{definition}

\subsection{Data Stream Model and Sliding Window}

\begin{definition}[Transaction Data Stream]
\label{def:stream}
A \emph{transaction data stream} $\mathit{TDS} = T_1, T_2, \ldots$ is
a potentially unbounded sequence of transactions arriving in real time.
A \emph{transaction-sensitive sliding window} $\SW$ of size $w$ at
time $N$ is the subsequence
$[T_{N-w+1}, T_{N-w+2}, \ldots, T_N]$.
As each new transaction $T_{N+1}$ arrives, $T_{N-w+1}$ is evicted.
\end{definition}

\subsection{Weighted Association Rules}

Following Ouyang and Huang~\cite{ouyang2015}, each item
$x \in \mathcal{I}$ is assigned a \emph{clinical weight}
$\omega(x) \in (0, 1]$ reflecting its significance (e.g., toxicity
index, interaction severity).

\begin{definition}[Weighted Support]
\label{def:wsup}
For an itemset $X = \{x_1, x_2, \ldots, x_k\}$ with
$\omega(x_1) \geq \omega(x_2) \geq \cdots \geq \omega(x_k)$,
the \emph{weighted support} in window $\SW$ is
\[
  \wsup(X)^{\SW}
    = \omega(X) \cdot \mathit{sup}(X)^{\SW}
    = \max_{i}\{\omega(x_i)\} \cdot \mathit{sup}(X)^{\SW}.
\]
\end{definition}

\begin{property}[Weighted Downward Closure~\cite{ouyang2015}]
\label{prop:closure}
If $X$ is a weighted frequent itemset, then every non-empty
subset $Y \subsetneq X$ is also weighted frequent.
\end{property}

\subsection{Regularity}

\begin{definition}[Regularity~\cite{tanbeer2010stream}]
\label{def:regularity}
Let $\mathit{tid}(X) = \{t_{j_1}, t_{j_2}, \ldots, t_{j_r}\}$ be the
ordered list of transaction identifiers in which itemset $X$ appears.
The \emph{regularity} of $X$ is the maximum inter-occurrence gap:
\[
  \mathit{reg}(X) = \max_{l=1}^{r-1}(t_{j_{l+1}} - t_{j_l}).
\]
$X$ is \emph{regular} if $\mathit{reg}(X) \leq \mathit{maxReg}$
for a user-defined threshold $\mathit{maxReg}$.
\end{definition}

% ─────────────────────────────────────────────────────────────
\section{The WNAR-SW Framework}
\label{sec:method}
% ─────────────────────────────────────────────────────────────

\subsection{Motivation}

Consider a chemical compound $c_A$ — a chemotherapy agent — with
clinical weight $\omega(c_A) = 0.95$.
Its raw frequency in a stream of 50{,}000 prescriptions is $0.008$
(400 occurrences), well below a typical $\minSup$ of $0.02$.
Standard Apriori and the algorithm of~\cite{sastry2025} would prune
$c_A$ before ever testing for negative associations.
Yet $c_A$ never co-occurs with a cardiac drug $c_B$
($\omega(c_B) = 0.90$), and this co-occurrence pattern — or rather,
its consistent \emph{absence} — is precisely the warning signal that
should reach the prescribing physician.

WNAR-SW resolves this through three novel definitions that jointly
ensure high-weight items receive proportionally relaxed frequency
thresholds and proportionally tightened regularity requirements.

\subsection{Weighted Regularity}

\begin{definition}[Weighted Regularity Threshold]
\label{def:weighted_regularity}
For an item $x$ with clinical weight $\omega(x)$, the
\emph{weighted regularity threshold} is
\[
  \maxRegW(x) = \mathit{maxReg} \cdot (1 - \omega(x)).
\]
Itemset $X$ is \emph{weighted-regular} if
$\mathit{reg}(X) \leq \min_{x \in X}\maxRegW(x)$.
\end{definition}

\begin{tcolorbox}[notebox]
\textbf{Intuition.}
A high-weight item (e.g., $\omega(x) = 0.95$) is subject to a very
tight regularity window ($\maxRegW = 0.05 \cdot \mathit{maxReg}$):
its appearance pattern must be \emph{highly consistent} before the
algorithm trusts it enough to retain it.
A low-weight item ($\omega(x) = 0.10$) inherits a loose threshold
($\maxRegW = 0.90 \cdot \mathit{maxReg}$) because mis-classifying its
pattern carries little clinical consequence.
\end{tcolorbox}

\subsection{Weighted Negative Support}

\begin{definition}[Weighted Negative Support]
\label{def:wnegsup}
For itemsets $A$ and $B$ satisfying
$\mathit{tid}(A) \cap \mathit{tid}(B) = \emptyset$ (complete mutual
exclusion in window $\SW$), the \emph{weighted negative support} is
\begin{equation}
  \wNegSup(A \Rightarrow \lnot B)^{\SW}
    = \omega(A)\cdot\omega(B)
      \cdot\bigl[\mathit{sup}(A)^{\SW} - \mathit{sup}(A \cup B)^{\SW}\bigr].
  \label{eq:wnegsup}
\end{equation}
When $\mathit{tid}(A) \cap \mathit{tid}(B) = \emptyset$, the term
$\mathit{sup}(A \cup B)^{\SW} = 0$, simplifying to:
\[
  \wNegSup(A \Rightarrow \lnot B)^{\SW}
    = \omega(A)\cdot\omega(B)\cdot\mathit{sup}(A)^{\SW}.
\]
\end{definition}

The product $\omega(A)\cdot\omega(B)$ acts as a severity multiplier:
two highly toxic incompatible chemicals receive a higher weighted
negative support than two benign ones, even at identical raw
frequencies.

\subsection{Weighted Negative Confidence}

\begin{definition}[Weighted Negative Confidence]
\label{def:wnegconf}
The \emph{weighted negative confidence} of rule
$A \Rightarrow \lnot B$ in window $\SW$ is
\begin{equation}
  \wNegConf(A \Rightarrow \lnot B)^{\SW}
    = \frac{\wNegSup(A \Rightarrow \lnot B)^{\SW}}{\wsup(A)^{\SW}}.
  \label{eq:wnegconf}
\end{equation}
\end{definition}

A rule is output to the alert system if and only if:
\[
  \wNegSup(A \Rightarrow \lnot B)^{\SW} \geq \mathit{minWNegSup}
  \quad\text{and}\quad
  \wNegConf(A \Rightarrow \lnot B)^{\SW} \geq \mathit{minWNegConf}.
\]

\subsection{Severity-Ranked Alert Score}

\begin{definition}[Alert Severity Score]
\label{def:severity}
The \emph{severity score} of a confirmed negative rule
$A \Rightarrow \lnot B$ is
\begin{equation}
  \mathit{sev}(A \Rightarrow \lnot B)
    = \wNegConf(A \Rightarrow \lnot B)^{\SW}
      \cdot \frac{\omega(A) + \omega(B)}{2}.
  \label{eq:severity}
\end{equation}
Alerts are sorted in descending order of $\mathit{sev}(\cdot)$
before delivery to the clinical decision-support interface.
\end{definition}

\subsection{Graceful Degradation}

\begin{proposition}[Backward Compatibility]
\label{prop:degradation}
\begin{enumerate}[label=(\roman*)]
  \item Setting $\omega(x) = 1$ for all $x \in \mathcal{I}$ and
        adopting flat regularity reduces WNAR-SW to the algorithm
        of Jammalamadaka and Budaraju~\cite{sastry2025}.
  \item Setting $\mathit{maxReg} = \infty$ (disabling regularity)
        and restricting to negative rules reduces WNAR-SW to a
        streaming weighted negative association miner extending
        Ouyang and Huang~\cite{ouyang2015}.
\end{enumerate}
\end{proposition}

This compatibility enables controlled benchmarking of the individual
contributions of weighting, regularity, and their interaction.

% ─────────────────────────────────────────────────────────────
\section{The WNAR-SW Algorithm}
\label{sec:algorithm}
% ─────────────────────────────────────────────────────────────

WNAR-SW operates as three concurrent processes sharing an
\emph{inverted table} data structure $\mathcal{V}$, which maps each
chemical code $c$ to the tuple
$(\mathit{tid\text{-}list}(c),\, \omega(c),\, \mathit{reg}(c))$.

\subsection{Process A — Historical Initialisation}

\begin{algorithm}[t]
\caption{Process-A: Historical Database Initialisation}
\label{alg:processA}
\small
\KwIn{Historical database $\mathcal{H}$; pharmaceutical weight table
  $\Omega$; thresholds $\minSup$, $\mathit{maxReg}$}
\KwOut{Initial weighted inverted table $\mathcal{V}$}

Read all records from $\mathcal{H}$ into memory\;
\ForEach{chemical $c$ appearing in $\mathcal{H}$}{
  $\omega(c) \leftarrow \Omega[c]$\;
  Compute $\mathit{tid\text{-}list}(c)$, $\mathit{sup}(c)$,
          $\mathit{reg}(c)$\;
  $\maxRegW(c) \leftarrow \mathit{maxReg} \cdot (1 - \omega(c))$\;
  \If{$\wsup(c) \geq \minSup$ \textbf{and}
      $\mathit{reg}(c) \leq \maxRegW(c)$}{
    Insert $(c, \mathit{tid\text{-}list}(c), \omega(c), \mathit{reg}(c))$
    into $\mathcal{V}$\;
  }
}
\Return $\mathcal{V}$\;
\end{algorithm}

\subsection{Process B — Sliding Window Stream Update}

\begin{algorithm}[t]
\caption{Process-B: Stream Ingestion and Inverted Table Update}
\label{alg:processB}
\small
\KwIn{Transaction data stream $\mathit{TDS}$; window size $w$;
  inverted table $\mathcal{V}$; thresholds}
\KwOut{Continuously updated and pruned $\mathcal{V}$}

Initialise sliding window buffer $\mathit{BUF}$ of capacity $w$\;
\While{\upshape stream $\mathit{TDS}$ is active}{
  \ForEach{incoming transaction $T_i \in \mathit{TDS}$}{
    \If{$|\mathit{BUF}| = w$}{
      Evict oldest transaction $T_{i-w}$ from $\mathit{BUF}$\;
      Remove $T_{i-w}.\mathit{tid}$ from all tid-lists in $\mathcal{V}$\;
    }
    Append $T_i$ to $\mathit{BUF}$\;
    \ForEach{chemical $c \in T_i$}{
      \eIf{$c \in \mathcal{V}$}{
        Append $T_i.\mathit{tid}$ to $\mathcal{V}[c].\mathit{tid\text{-}list}$\;
        Update $\mathit{sup}(c)$, $\mathit{reg}(c)$\;
      }{
        $\omega(c) \leftarrow \Omega[c]$\;
        Insert new entry for $c$ into $\mathcal{V}$\;
      }
    }
  }
  \tcp{Prune after each window slide}
  \ForEach{entry $c \in \mathcal{V}$}{
    \If{$\wsup(c)^{\SW} < \minSup$ \textbf{or}
        $\mathit{reg}(c) > \maxRegW(c)$}{
      Remove $c$ from $\mathcal{V}$\;
    }
  }
}
\end{algorithm}

\subsection{Process C — Weighted Negative Association Detection}

\begin{algorithm}[t]
\caption{Process-C: Negative Association Mining and Alert Generation}
\label{alg:processC}
\small
\KwIn{Pruned inverted table $\mathcal{V}$; thresholds
  $\mathit{minWNegSup}$, $\mathit{minWNegConf}$; drug-chemical
  mapping $\mathcal{M}$}
\KwOut{Ranked alert set $\mathcal{A}$}

$\mathcal{A} \leftarrow \emptyset$\;
\ForEach{pair $(A, B)$ in $\mathcal{V}$ with $A \neq B$}{
  \If{$\mathcal{V}[A].\mathit{tid\text{-}list} \cap
      \mathcal{V}[B].\mathit{tid\text{-}list} = \emptyset$}{
    Compute $\wNegSup(A \Rightarrow \lnot B)^{\SW}$
    using~\eqref{eq:wnegsup}\;
    Compute $\wNegConf(A \Rightarrow \lnot B)^{\SW}$
    using~\eqref{eq:wnegconf}\;
    \If{$\wNegSup \geq \mathit{minWNegSup}$ \textbf{and}
        $\wNegConf \geq \mathit{minWNegConf}$}{
      $s \leftarrow$ severity score from~\eqref{eq:severity}\;
      Map $A, B$ to drug names via $\mathcal{M}$\;
      Insert alert $(A \Rightarrow \lnot B, s)$ into $\mathcal{A}$\;
    }
  }
}
Sort $\mathcal{A}$ by severity score $s$ descending\;
\Return $\mathcal{A}$\;
\end{algorithm}

\subsection{Data Structure: Weighted Inverted Table}

The inverted table $\mathcal{V}$ stores, for each retained chemical
$c$:
\[
  \mathcal{V}[c] = \bigl(
    \mathit{tid\text{-}list}(c),\;
    \omega(c),\;
    \mathit{sup}(c)^{\SW},\;
    \mathit{reg}(c),\;
    \wsup(c)^{\SW},\;
    \maxRegW(c)
  \bigr).
\]
This representation allows Process C to test mutual exclusion in
$O(|\mathit{tid\text{-}list}(A)| + |\mathit{tid\text{-}list}(B)|)$
time via sorted-list intersection, and to compute all weighted
measures in $O(1)$ per pair.

% ─────────────────────────────────────────────────────────────
\section{Theoretical Analysis}
\label{sec:analysis}
% ─────────────────────────────────────────────────────────────

\subsection{Correctness of Weighted Anti-Monotonicity Under Negative Mining}

\begin{theorem}[Weighted Closure Under Negative Extension]
\label{thm:closure}
Let $A$ and $B$ be disjoint itemsets such that
$A \Rightarrow \lnot B$ is a weighted negative association rule
under WNAR-SW.
Then for any $A' \subseteq A$ and $B' \subseteq B$ with
$A' \neq \emptyset$ and $B' \neq \emptyset$:
\[
  \wNegSup(A' \Rightarrow \lnot B')^{\SW}
    \geq \wNegSup(A \Rightarrow \lnot B)^{\SW}.
\]
\end{theorem}

\begin{proof}
By \cref{prop:closure} (weighted downward closure),
$\wsup(A')^{\SW} \geq \wsup(A)^{\SW}$ and
$\wsup(B')^{\SW} \geq \wsup(B)^{\SW}$.
Since $\mathit{tid}(A) \cap \mathit{tid}(B) = \emptyset$ implies
$\mathit{tid}(A') \cap \mathit{tid}(B') = \emptyset$
(as $\mathit{tid}(A') \supseteq \mathit{tid}(A)$ and
$\mathit{tid}(B') \supseteq \mathit{tid}(B)$), we have
$\mathit{sup}(A' \cup B')^{\SW} = 0$.
Thus $\mathit{sup}(A')^{\SW} - \mathit{sup}(A' \cup B')^{\SW}
      = \mathit{sup}(A')^{\SW} \geq \mathit{sup}(A)^{\SW}$.
Moreover, $\omega(A') \geq \omega(A)$ and $\omega(B') \geq \omega(B)$
by the max-weight construction of \cref{def:wsup}.
Therefore $\wNegSup(A' \Rightarrow \lnot B')^{\SW}
  = \omega(A')\cdot\omega(B')\cdot\mathit{sup}(A')^{\SW}
  \geq \omega(A)\cdot\omega(B)\cdot\mathit{sup}(A)^{\SW}
  = \wNegSup(A \Rightarrow \lnot B)^{\SW}$.
\end{proof}

\begin{remark}
\Cref{thm:closure} guarantees that Process C never misses a weighted
negative association by focusing on minimal antecedent/consequent
pairs: if the pair $(A, B)$ fails the threshold, so will any
$(A'', B'')$ with $A'' \supsetneq A$ or $B'' \supsetneq B$.
This property is exploited during candidate pruning.
\end{remark}

\subsection{Complexity Analysis}

\begin{proposition}[Space Complexity]
WNAR-SW maintains $O(|\mathcal{V}| \cdot w)$ space, where
$|\mathcal{V}|$ is the number of retained chemicals (bounded by
$|\mathcal{I}|$) and $w$ is the window size.
This is a fixed upper bound independent of stream length.
\end{proposition}

\begin{proposition}[Per-Transaction Time Complexity]
Each incoming transaction $T_i$ of size $|T_i|$ incurs
$O(|T_i| \cdot \log|\mathcal{V}|)$ cost for inverted table lookup
and update, plus $O(|\mathcal{V}|)$ for the pruning pass.
Process C runs asynchronously; its per-window cost is
$O(|\mathcal{V}|^2 \cdot w)$ for exhaustive pairwise intersection.
\end{proposition}

% ─────────────────────────────────────────────────────────────
\section{Discussion}
\label{sec:discussion}
% ─────────────────────────────────────────────────────────────

\subsection{Worked Example}
\label{subsec:example}

We illustrate WNAR-SW on a synthetic sliding window of $w = 30$
transactions drawn from a medical dataset.
After Process B, the pruned inverted table $\mathcal{V}$ retains two
chemicals of interest:

\smallskip
\begin{center}
\small
\begin{tabular}{@{}lccc@{}}
  \toprule
  Chemical & $\omega$ & $\mathit{sup}^{\SW}$ & $\wsup^{\SW}$ \\
  \midrule
  $c_A$ (chemo agent) & 0.95 & $4/30 \approx 0.133$ & $0.126$ \\
  $c_B$ (cardiac drug) & 0.90 & $5/30 \approx 0.167$ & $0.150$ \\
  \bottomrule
\end{tabular}
\end{center}
\smallskip

\noindent
Since $\mathit{tid}(c_A) \cap \mathit{tid}(c_B) = \emptyset$:
\begin{align*}
  \wNegSup(c_A \Rightarrow \lnot c_B)^{\SW}
    &= 0.95 \times 0.90 \times 0.133 \approx 0.114,\\
  \wNegConf(c_A \Rightarrow \lnot c_B)^{\SW}
    &= 0.114 / 0.126 \approx 0.905,\\
  \mathit{sev}(c_A \Rightarrow \lnot c_B)
    &= 0.905 \times \tfrac{0.95 + 0.90}{2} \approx 0.837.
\end{align*}

With $\mathit{minWNegSup} = 0.10$ and $\mathit{minWNegConf} = 0.80$,
both thresholds are satisfied and this rule is output as a
\textbf{HIGH} severity alert ($\mathit{sev} = 0.837$).
Under~\cite{sastry2025}, $c_A$ would have been pruned at the frequency
threshold because $\mathit{sup}(c_A) = 0.133 < \minSup = 0.20$
(a typical setting), and the interaction would go undetected.

\subsection{Expected Experimental Behaviour}

Based on the results reported in~\cite{sastry2025} and the
mathematical structure of WNAR-SW, we project the following
behavioural characteristics:

\begin{enumerate}[leftmargin=*, itemsep=3pt]

  \item \textbf{Fewer false-positive alerts than~\cite{sastry2025}.}
        Weighted regularity filters spurious patterns from high-weight
        items more aggressively; the severity ranking further ensures
        only clinically significant interactions reach the physician.

  \item \textbf{Detection of interactions missed by~\cite{sastry2025}.}
        High-weight chemicals with raw frequency below $\minSup$ are
        retained under the weighted support criterion, enabling
        detection of rare but dangerous incompatibilities.

  \item \textbf{Ranked output quality.}
        The severity score orders alerts so that interactions between
        the most dangerous chemicals appear first, enabling triage in
        high-volume ICU settings.

  \item \textbf{Threshold sensitivity.}
        Increasing $\mathit{minWNegSup}$ reduces alert volume rapidly;
        decreasing it recovers marginal interactions.
        A calibration study on real prescription data is required to
        determine optimal settings per clinical specialty.

\end{enumerate}

\subsection{Weight Assignment in Practice}

WNAR-SW is agnostic as to the source of weights $\omega(\cdot)$.
We recommend the following hierarchy:
(i) normalised adverse-effect severity scores from pharmacological
databases (e.g., DrugBank, SIDER);
(ii) expert-elicited weights from a clinical pharmacology panel;
(iii) uniform weights ($\omega \equiv 1$), which reproduces
\cite{sastry2025} and provides a controlled baseline.
Machine-learned weight estimation from historical adverse-event records
is a natural direction for future work.

\subsection{Limitations}

\begin{enumerate}[leftmargin=*, label=\textbf{Lim.\arabic*.}]
  \item \textbf{Pairwise scope.}
        WNAR-SW detects two-chemical negative associations.
        Three-body interactions (\eg $A \Rightarrow \lnot (B \cup C)$)
        are not currently handled; extension to weighted negative
        $k$-itemsets is deferred to future work.

  \item \textbf{Synthetic validation only.}
        Initial experiments use synthetic data generated to protect
        patient privacy.
        Validation on anonymised clinical prescription streams, subject
        to ethics board approval, is required before clinical deployment.

  \item \textbf{Threshold proliferation.}
        WNAR-SW introduces $\mathit{minWNegSup}$, $\mathit{minWNegConf}$,
        and $\mathit{maxReg}$ in addition to the thresholds of prior work.
        A sensitivity analysis and automated threshold recommendation
        component are planned.

  \item \textbf{Concept drift.}
        Prescribing patterns may shift systematically over time
        (\eg during a pandemic).
        An adaptive threshold mechanism responsive to drift is a priority
        extension.
\end{enumerate}

% ─────────────────────────────────────────────────────────────
\section{Conclusion}
\label{sec:conclusion}
% ─────────────────────────────────────────────────────────────

This paper has identified and formally characterised the gap between
two state-of-the-art extensions of Apriori: the frequency-regularity
negative association miner of Jammalamadaka and Budaraju~\cite{sastry2025}
and the weighted association rule miner of Ouyang and Huang~\cite{ouyang2015}.
Neither algorithm — nor any prior work in the reviewed literature —
simultaneously handles data streams, negative associations, item
weighting, and temporal regularity.

We proposed WNAR-SW, which closes this gap through three unified
definitions: weighted regularity (\cref{def:weighted_regularity}),
weighted negative support (\cref{def:wnegsup}), and weighted negative
confidence (\cref{def:wnegconf}).
We proved the correctness of the weighted anti-monotonicity property
under the negative mining extension (\cref{thm:closure}) and
demonstrated that the algorithm degrades gracefully to either baseline
under appropriate parameter settings (\cref{prop:degradation}).
A worked example confirmed that WNAR-SW detects a clinically critical
drug incompatibility that both prior algorithms would silently miss.

The clinical impact of this contribution extends beyond algorithmic
novelty.
Real-time, severity-ranked drug interaction alerts — anchored to both
pharmacological weights and temporal regularity — represent a step
toward the kind of automated prescription oversight that could have
prevented patient harm during the COVID-19 pandemic and can continue
to do so in everyday clinical practice.

Future work will focus on:
extension to weighted negative $k$-itemsets for multi-drug
interactions;
machine-learned weight estimation;
adaptive drift-tolerant thresholding;
and validation on de-identified real-world prescription streams in
collaboration with clinical partners.

% ─────────────────────────────────────────────────────────────
%  ACKNOWLEDGEMENTS
% ─────────────────────────────────────────────────────────────
\section*{Acknowledgements}
The authors thank [Supervisor Name] for guidance and the Department of
Computer Science at [University Name] for computational resources.

% ─────────────────────────────────────────────────────────────
%  REFERENCES
% ─────────────────────────────────────────────────────────────
\bibliographystyle{unsrt}
\begin{thebibliography}{99}

\bibitem{agrawal1993}
R.~Agrawal, T.~Imielinski, and A.~Swami,
``Mining association rules between sets of items in large databases,''
in \textit{Proc. ACM SIGMOD}, pp.~207--216, 1993.

\bibitem{agrawal1994}
R.~Agrawal and R.~Srikant,
``Fast algorithms for mining association rules,''
in \textit{Proc. 20th VLDB Conference}, pp.~487--499, 1994.

\bibitem{sastry2025}
S.~K.~R.~Jammalamadaka and R.~R.~Budaraju,
``Finding negative associations from medical data streams based on
frequent and regular patterns,''
\textit{Contemporary Mathematics}, vol.~6, no.~2, pp.~1434--1454, 2025.
\textsc{doi}: \href{https://doi.org/10.37256/cm.6220256229}%
{10.37256/cm.6220256229}.

\bibitem{ouyang2015}
W.~Ouyang and Q.~Huang,
``Mining weighted association rules in data streams with a sliding
window,''
in \textit{Proc. Int.\ Conf.\ on Computer Science and Intelligent
Communication (CSIC)}, pp.~271--274, 2015.

\bibitem{ouyang2013}
W.~Ouyang,
``Mining positive and negative association rules in data streams with
a sliding window,''
in \textit{Proc. 4th Global Congress on Intelligent Systems}, 2013.

\bibitem{zaki2003}
M.~J.~Zaki,
``Fast vertical mining using diffsets,''
in \textit{Proc. 9th ACM SIGKDD}, pp.~326--335, 2003.

\bibitem{savasere1998}
A.~Savasere, E.~Omiecinski, and S.~Navathe,
``Mining for strong negative associations in a large database of
customer transactions,''
in \textit{Proc. 14th ICDE}, pp.~494--502, 1998.

\bibitem{antonie2004}
M.~Antonie and O.~R.~Zaiane,
``Mining positive and negative association rules: An approach for
confined rules,''
in \textit{Proc. 8th PKDD}, pp.~27--38, 2004.

\bibitem{cornelis2006}
C.~Cornelis, P.~Yan, X.~Zhang, and G.~Chen,
``Mining positive and negative association rules from large databases,''
in \textit{Proc. IEEE Conf.\ on Cybernetics and Intelligent Systems},
pp.~1--6, 2006.

\bibitem{dong2006}
X.~Dong, F.~Sun, X.~Han, and R.~Hou,
``Study of positive and negative association rules based on
multi-confidence and chi-squared test,''
in \textit{Proc. ADMA}, pp.~100--109, 2006.

\bibitem{lin2006}
C.-H.~Lin, D.-Y.~Chiu, Y.-H.~Wu, and A.~L.-P.~Chen,
``Mining frequent itemsets from data streams with a time-sensitive
sliding window,''
in \textit{Proc. SIAM SDM}, 2006.

\bibitem{tanbeer2010stream}
S.~K.~Tanbeer, C.~F.~Ahmed, and B.-S.~Jeong,
``Mining regular patterns in data streams,''
in \textit{Proc. DASFAA}, pp.~399--413, 2010.

\bibitem{tanbeer2009sliding}
S.~K.~Tanbeer, C.~F.~Ahmed, B.-S.~Jeong, and Y.-K.~Lee,
``Sliding window-based frequent pattern mining over data streams,''
\textit{Information Sciences}, vol.~179, no.~22, pp.~3843--3865, 2009.

\bibitem{cai1998}
C.~H.~Cai, A.~W.-C.~Fu, C.-H.~Cheng, and W.~W.~Kwong,
``Mining association rules with weighted items,''
in \textit{Proc. IDEAS}, pp.~68--77, 1998.

\bibitem{bagui2019}
S.~Bagui and P.~C.~Dhar,
``Positive and negative association rule mining in Hadoop's MapReduce
environment,''
\textit{Journal of Big Data}, vol.~6, p.~75, 2019.

\bibitem{wei2020}
J.~Wei, Z.~Lu, K.~Qiu, P.~Li, and A.~H.~Sun,
``Predicting drug risk level from adverse drug reactions using SMOTE
and machine learning approaches,''
\textit{IEEE Access}, vol.~8, pp.~185761--185775, 2020.

\bibitem{zhang2020}
J.~Zhang, W.~Liu, and P.~Wang,
``Drug-drug interaction extraction from Chinese biomedical literature
using distant supervision,''
in \textit{Proc. IEEE ICKG}, pp.~593--598, 2020.

\end{thebibliography}

\end{document}
